% generated by the QuadraSHAP experiments -- do not edit by hand
\begin{table}[t]
\centering
\small
\caption{The a priori certificate over a population of instances at fixed total relative variation. For each feature count, the RBF scale is calibrated so the population mean is $\bar\Lambda=40$. For every instance and every node count below the exactness threshold, the certified bound $B$ is computed from the factor tables before any quadrature runs and compared with the observed error $E$ against the exact rule. A point is \emph{informative} when the certificate is still above the rounding level of the attributions; coverage is the fraction of $(\text{instance}, m_q)$ pairs with $E \le B$ up to that rounding level, and the slack columns are taken over the points whose observed error is above it.}
\label{tab:certificate-population}
\begin{tabular}{@{}lrrrrrrrrrr@{}}
\toprule
$d$ & instances & $\gamma d$ & $\bar\Lambda$ & sd$(\Lambda)$ & median $m_q^\star$ & points & coverage & median $\log_{10}\frac{B}{E}$ & 90th pct. & Spearman $(B,E)$ \\
\midrule
100 & 60 & 29.15 & 40.0 & 2.1 & 12 & 540 & 1.0000 & 2.8 & 4.5 & 0.961 \\
500 & 40 & 19.41 & 40.0 & 1.2 & 12 & 360 & 1.0000 & 2.8 & 3.2 & 0.939 \\
1000 & 25 & 18.97 & 40.0 & 0.9 & 12 & 225 & 0.6667 & 2.5 & 3.0 & 0.933 \\
\bottomrule
\end{tabular}
\\[2pt] \footnotesize Kernel ridge regression with an isotropic RBF kernel, $n=300$ training points, neutral-factor value function; the same node grid is used for every $d$.
\end{table}
